\section{Urbilateria to early chordates: the material substrate of agency}\label{04_urbilateria-urbilateria-to-early-chordates-the-material-substrate-of-agency}

The path from the last common bilaterian ancestor to the first vertebrates assembled, over roughly 60--80 million years of Ediacaran-Cambrian evolution, the cellular machinery that still implements goal-directed behaviour in human nervous systems. \textbf{The story is one of co-option, not invention}: every classical neurotransmitter is a repurposed metabolic intermediate, every synaptic protein has homologs in unicellular eukaryotes, and the modern bilaterian synapse is a polarised re-organisation of secretory and signalling modules that long predate neurons. What was added during this window was not new chemistry but new \textbf{architecture} --- bilateral bodies forced to govern their movement through structured environments, the assembly of a centralised nerve cord (perhaps once, perhaps more than once), and the chordate experiment with a dorsal hollow tube whose deepest molecular partitions are still recognisable in the lamprey forebrain. Folk-psychological terms like ``decision,'' ``purpose,'' and ``intention'' attach to this scaffolding only loosely; the stable referent is the substrate itself, and the substrate is, materially, an evolved chemistry that can break.

This chapter traces that scaffolding from the urbilaterian --- already in possession of Hox clusters, a Six3/Otx/Gbx/Hox axial code, a BMP-Chordin dorsoventral cassette, neuropeptidergic GPCR networks, and the full small-molecule transmitter repertoire --- through deuterostome divergence, the chordate body plan, and finally to the lamprey, whose half-billion-year-old brain already implements the action-selection circuitry that later vertebrate evolution would elaborate but not invent.

\subsection{The governance problem of bilateral bodies}\label{04_urbilateria-the-governance-problem-of-bilateral-bodies}

Cnidarians and ctenophores are predominantly radial: their environments are isotropic, prey may approach from any direction, and a diffuse nerve net suffices to coordinate isotropic responses (Arendt, Tosches \& Marlow 2016, \emph{Nat Rev Neurosci}). A bilateral body with a polarised anteroposterior axis and a through-gut faces a categorically different problem. One end consistently encounters novel stimuli before the rest. Mouth and anus are spatially separated, requiring continuous peristaltic coordination of intake, processing, and egestion. The body moves directionally, so the cost of mis-prediction at the leading end falls on the whole organism. Trestman (2013, \emph{Biological Theory}) formalised this as the problem of \textbf{complex active bodies}: only three bilaterian phyla --- arthropods, cephalopod molluscs, and chordates --- independently evolved highly cephalised, complex active sensorimotor lifestyles, \href{https://en.wikipedia.org/wiki/Cephalization}{Wikipedia} and the architecture they converged on shares low-latency anterior sensory clusters, hierarchical integration, and central pattern generators feeding paired motor outputs.

Cephalisation is the convergent solution but not a unitary trait. Hox1--5 genes pattern cephalic structures in vertebrates and arthropods, but Fernández-Nicolás and Campuzano (2021) argue these Hox cephalic functions were acquired \textbf{independently} in the two lineages \href{https://pmc.ncbi.nlm.nih.gov/articles/PMC8374599/}{PubMed Central} --- a parallel co-option of the same regulatory toolkit rather than direct homology of the head. Finnerty (2005, \emph{BioEssays}) provides a useful corrective: bilaterality may have arisen first for \textbf{internal} transport demands --- efficient circulation through an elongated body --- with directed locomotion and cephalisation as later consequences rather than original causes. Importantly, cephalisation is not inevitable; many bilaterian lineages, especially among echinoderms, have lost it secondarily.

The selection pressures inferred from comparative work are more concrete than the broad governance language suggests: priority of anterior sensory encounter during forward locomotion, reduction of conduction delays between sensors and effectors, multimodal integration to support behavioural choice, and the Cambrian establishment of multi-tiered predator-prey food webs (Erwin et al.~2011, \emph{Science}; Peterson et al.~2008, \emph{Phil Trans R Soc B}). What had to be solved, materially, was \textbf{action selection} --- given multiple available motor patterns and multimodal sensory input, suppress all but one. Strausfeld and Hirth (2013, \emph{Science}) argue the deep homology of the arthropod central complex and the vertebrate basal ganglia \href{https://evodevojournal.biomedcentral.com/articles/10.1186/2041-9139-4-27}{BioMed Central} reflects exactly this constraint, \href{https://www.science.org/doi/abs/10.1126/science.1231828}{Science} \href{https://scholar.google.com/scholar_lookup?title=Deep+homology+of+arthropod+central+complex+and+vertebrate+basal+ganglia&amp=&author=N.+J.+Strausfeld&amp=&author=F.+Hirth&amp=&publication_year=2013&amp=&journal=Science&amp=&pages=157-161&amp=&doi=10.1126/science.1231828&amp=&pmid=23580521}{Google Scholar} conserved or convergently re-evolved across 600 million years.

\subsection{Neurotransmitters as repurposed metabolites}\label{04_urbilateria-neurotransmitters-as-repurposed-metabolites}

Every classical neurotransmitter is a metabolite first and a signal second. This is not a rhetorical framing but a consequence of biochemistry: the molecules used in synaptic transmission are pre-existing intermediates of central carbon, nitrogen, and aromatic-amino-acid metabolism, and the receptors that detect them descend from far older ion channels and seven-transmembrane receptors used by bacteria, archaea, plants, fungi, and unicellular holozoans. Neurons co-opted what was already chemically present and metabolically essential.

\textbf{Calcium} is the foundational case. Its role as universal intracellular signal was forced by chemistry, not designed: calcium phosphate is essentially insoluble (Kₛₚ for Ca₃(PO₄)₂ ≈ 2 × 10⁻²⁹), and life is built on phosphate. Every cell that uses ATP must keep cytosolic Ca²⁺ around 100 nM against a 1--2 mM extracellular concentration, or it precipitates its own metabolism. Williams (1999) and Carafoli (2003, \emph{Trends Biochem Sci}; 2005, \emph{FEBS J}) argue this exclusion gradient was a chemical necessity that signalling repurposed. The ``calcium toolkit'' of pumps (PMCA, SERCA), exchangers (NCX), channels (Cav, ryanodine receptors, IP₃Rs), and EF-hand sensors (calmodulin, calcineurin) is conserved across all eukaryotes and detectable in prokaryotes (Berridge, Bootman \& Roderick 2003, \emph{Nat Rev Mol Cell Biol}; Cai 2017). Voltage-gated Cav, Nav, and Kv channels descend from a prokaryotic six-transmembrane K⁺ channel ancestor; \href{https://pmc.ncbi.nlm.nih.gov/articles/PMC12862854/}{PubMed Central} \href{https://journals.biologists.com/jeb/article/218/4/515/14125/Evolution-of-voltage-gated-ion-channels-at-the}{The Company of Biologists} the four-domain pseudotetrameric Cav is older than Nav, \href{https://www.nature.com/articles/s41598-018-21897-7}{Nature} and Liebeskind, Hillis and Zakon (2011, \emph{PNAS}) showed that \textbf{Na⁺-selective channels were derived independently in animals and bacteria from Ca²⁺-selective ancestors}. \href{https://pmc.ncbi.nlm.nih.gov/articles/PMC7723092/}{PubMed Central} Decoupling rapid electrical excitability from intracellular Ca²⁺ signalling was a key animal innovation \href{https://pmc.ncbi.nlm.nih.gov/articles/PMC7723092/}{PubMed Central} enabling fast action potentials without flooding the cytosol with the universal second messenger.

\textbf{Glutamate} sits at the intersection of carbon and nitrogen metabolism --- the α-ketoglutarate/glutamate interconversion is the principal nitrogen entry point for plants and microbes, and free cytosolic glutamate concentrations in neurons reach 1--10 mM. Ionotropic glutamate receptors \href{https://pmc.ncbi.nlm.nih.gov/articles/PMC9233959/}{PubMed Central} are a modular fusion of a periplasmic-binding-protein-like \href{https://pmc.ncbi.nlm.nih.gov/articles/PMC9233959/}{PubMed Central} Venus flytrap ligand-binding domain and an inverted bacterial KcsA-type K⁺ channel pore (Chen et al.~1999 \href{https://pmc.ncbi.nlm.nih.gov/articles/PMC9233959/}{PubMed Central} on cyanobacterial GluR0; Kuner, Seeburg \& Guy 2003). \href{https://www.ncbi.nlm.nih.gov/pmc/articles/PMC2924276/}{nih} Plants encode 20 GLR genes (Lam et al.~1998, \emph{Nature}); \href{https://physoc.onlinelibrary.wiley.com/doi/10.1113/JP279026}{The Physiological Society} they function as amino-acid-gated Ca²⁺ channels \href{https://www.ncbi.nlm.nih.gov/pmc/articles/PMC2924276/}{nih} in defence, pollen-tube growth and the long-distance wound Ca²⁺ waves of Toyota et al.~(2018, \emph{Science}). \textbf{Plant GLRs are no closer to metazoan iGluRs than to bacterial ones}, indicating divergence near the eukaryote stem. \href{https://pmc.ncbi.nlm.nih.gov/articles/PMC3483616/}{PubMed Central} \href{https://www.frontiersin.org/articles/10.3389/fpls.2012.00235/full}{Frontiers} Ramos-Vicente \href{https://www.researchgate.net/scientific-contributions/David-Ramos-Vicente-2129164358}{ResearchGate} et al.~(2018, \emph{eLife}) reorganised metazoan iGluRs into \href{https://www.biorxiv.org/content/10.1101/2024.06.25.600656v1.full}{bioRxiv} AKDF, NMDA, Epsilon and a sponge-only Lambda subfamily. Crystallography of ctenophore (Alberstein et al.~2015, \emph{PNAS}) and \emph{Trichoplax} iGluRs shows \href{https://www.nature.com/articles/s42003-025-08402-3}{Nature} endogenous \textbf{glycine} binding rather than glutamate --- ligand specificity has switched repeatedly. \href{https://www.nature.com/articles/s42003-025-08402-3}{Nature} Metabotropic glutamate receptors, GABA-B, the Ca²⁺-sensing receptor and T1R taste receptors are all class C GPCRs sharing the bilobed extracellular Venus flytrap from bacterial periplasmic amino-acid-binding proteins, and mGluR-like sequences are present in sponges \href{https://www.sciencedirect.com/topics/biochemistry-genetics-and-molecular-biology/ionotropic-glutamate-receptor}{ScienceDirect} and choanoflagellates (King et al.~2008; Krishnan et al.~2012).

\textbf{GABA} is glutamate decarboxylated by glutamate decarboxylase (GAD) using PLP. Its original role is acid-stress and metabolic shunting: in \emph{E. coli}, GAD consumes a cytoplasmic proton, allowing survival in stomach acid below pH 2.5. Plants use GABA in heat, anoxia, wound and pathogen responses, and it already gates ion channels in plants --- ALMT (aluminium-activated malate transporter) channels and GORK potassium channels (Ramesh et al.~2015; Bouché \& Fromm 2004). GABA-A and GABA-C are pentameric Cys-loop chloride channels; GABA-B is a class C GPCR. The Cys-loop superfamily descends from a bacterial pentameric ``Pro-loop'' channel --- Tasneem, Iyer, Jakobsson and Aravind (2005) identified prokaryotic homologs (GLIC from \emph{Gloeobacter}, ELIC from \emph{Erwinia}) --- with disulfide-bridged Cys-loop derivatives diversifying in eukaryotes into nAChRs, GABA-A, glycine, 5-HT₃ and invertebrate glutamate-gated chloride channels \href{https://pmc.ncbi.nlm.nih.gov/articles/PMC10480131/}{PubMed Central} (Ortells \& Lunt 1995; Jaiteh, Taly \& Hénin 2016).

\textbf{Dopamine} is two enzymatic steps from tyrosine: tyrosine hydroxylase (TH) yields L-DOPA, and aromatic L-amino-acid decarboxylase (AADC/DDC) decarboxylates it. \href{https://pmc.ncbi.nlm.nih.gov/articles/PMC3065393/}{PubMed Central} The three pterin-dependent aromatic amino acid hydroxylases --- phenylalanine, tyrosine, and tryptophan hydroxylase --- share a deep common ancestor with bacterial substrate-promiscuous enzymes (Fitzpatrick 1999; Mora-Soares et al.~2022). The original biochemical role was likely \textbf{antioxidant production} --- hydroxylated aromatics scavenge reactive oxygen species --- predating any signalling role. Dopamine reaches millimolar concentrations in plants like banana and \emph{Mucuna} as antioxidant; it is detectable in bacteria, fungi, sponges, cnidarians and all bilaterians. \emph{Drosophila} TH (encoded by \emph{pale}) has two splice isoforms, one for CNS, one for \textbf{cuticle sclerotisation} in epidermis \href{https://pubmed.ncbi.nlm.nih.gov/10358021/}{PubMed} --- a striking reminder that the gene serves multiple ancient roles and that synaptic dopaminergic signalling is one application among several. All five mammalian dopamine receptors are class A rhodopsin-family GPCRs.

\textbf{Serotonin} is similarly derived from tryptophan via tryptophan hydroxylase, and its presence in plants, algae, fungi, bacteria, sponges and hydra (Erland, Turi \& Saxena 2016) is consistent with an ancient role as antioxidant --- closely related to melatonin, which Manchester et al.~(2015) and Tan et al.~(2009) propose evolved when oxygenic photosynthesis began oxidising the biosphere. Most of the fourteen mammalian 5-HT receptors are class A GPCRs; the exception is \textbf{5-HT₃}, a Cys-loop pentameric cation channel co-opted from the nAChR-related lineage. In gastropod molluscs --- \emph{Aplysia}, \emph{Lymnaea}, \emph{Pleurobranchaea}, \emph{Tritonia}, \emph{Clione} --- serotonin from giant cerebral neurons modulates feeding, locomotion, arousal and learning, including Kandel's classic gill-withdrawal sensitisation. These ancient feeding/locomotor roles probably represent the ancestral bilaterian function of the serotonergic system.

\textbf{Acetylcholine} is in some respects the most mundane case: choline (from phospholipid turnover) plus acetyl-CoA (from central catabolism) makes ACh, and both substrates are universal. ACh has been present for \textasciitilde3 billion years in bacteria, blue-green algae, fungi, protozoa and plants, where it modulates water uptake, photosynthesis and growth (Wessler \& Kirkpatrick 2008, \emph{Br J Pharmacol}). Nicotinic AChRs are pentameric Cys-loop cation channels; muscarinic AChRs are class A GPCRs (M1--M5). The crystal structure of the snail \emph{Lymnaea} AChBP (Brejc et al.~2001, \emph{Nature}) first revealed the ancestral ligand-binding-site architecture of the entire Cys-loop family.

\textbf{Other monoamines} complete the picture. Octopamine and tyramine are the ``noradrenaline equivalents'' in protostomes --- but Bauknecht and Jékely (2017, \emph{BMC Biol}) showed that adrenergic α1, α2 receptor families plus octopamine α/β and tyramine 1/2 receptors all coexisted in the urbilaterian. Vertebrates \textbf{lost} octopamine/tyramine receptors; most insects and nematodes lost adrenergic ones. Norepinephrine and octopamine are not replacements but ancestral coexistents. Histamine is a histidine-decarboxylation product gating GPCRs in mammals and Cys-loop chloride channels in \emph{Drosophila} photoreceptors.

The vesicular machinery is similarly recycled. All vesicular neurotransmitter transporters --- VGLUT (SLC17), VMAT/VAChT (SLC18), VGAT (SLC32) --- are derived from the major-facilitator superfamily, energised by V-type H⁺-ATPases that themselves descend from ancestral eukaryotic vacuolar acidification machinery. The VAChT gene is nested within an intron of ChAT, a synteny conserved from invertebrates to vertebrates. Burkhardt (2015, \emph{J Exp Biol}) and Burkhardt and Jékely (2021) document SNARE/Munc18 vesicle-fusion machinery, complete neurosecretory exocytosis modules, voltage-gated Ca²⁺ channels and Homer/Shank/PSD-95 scaffolds in choanoflagellates \href{https://www.semanticscholar.org/paper/The-origin-and-evolution-of-synaptic-proteins-lead-Burkhardt/296eb4866f003b9b1ce6882845cc87f265606057}{Semantic Scholar +2} --- all the major synaptic components are pre-metazoan.

The synthesis is unambiguous: glutamate (TCA/N hub), GABA (acid-stress shunt), dopamine and serotonin (Tyr/Trp antioxidant pathways), acetylcholine (universal metabolic byproduct), histamine (His decarboxylation) --- every classical neurotransmitter is a metabolite first. The biosynthetic enzymes are co-opted promiscuous bacterial ancestors. The receptors descend from prokaryotic and eukaryotic ion channels and 7TM receptors. \textbf{Calcium provides the chemical universal that all signalling routes ultimately couple to.} When tyrosine hydroxylase fails in a human substantia nigra, Parkinsonism follows; the same biosynthetic step has worked, and failed, since well before bilaterians.

\subsection{The synapse as a gradual assembly}\label{04_urbilateria-the-synapse-as-a-gradual-assembly}

The chemical synapse is best understood not as a singular novelty but as the tightening of cellular machinery that long predated nervous systems into a polarised cell-cell signalling junction. Comparative proteomics identified roughly 1,400 proteins in the human postsynaptic density, of which a ``core protosynapse'' set is traceable to unicellular eukaryotes --- the average human-prokaryote alignment of PSD-domain homologs covers 96.5\% of conserved Pfam domains \href{https://www.frontiersin.org/journals/neuroscience/articles/10.3389/fnins.2011.00044/full}{Frontiers} (Emes \& Grant 2011, \emph{Front Neurosci}; Ryan \& Grant 2009, \emph{Nat Rev Neurosci}).

The \textbf{presynaptic} module assembled deep in unicellular ancestors. Burkhardt et al.~(2011, \emph{PNAS}) found that \emph{Monosiga brevicollis} --- a choanoflagellate --- already possesses \href{https://www.researchgate.net/publication/272496772_The_origin_and_evolution_of_synaptic_proteins_-_Choanoflagellates_lead_the_way}{ResearchGate} the Munc18/syntaxin-1 binary complex with vertebrate-like architecture and tight closed-syntaxin regulation, previously thought neuron-specific. \href{https://pmc.ncbi.nlm.nih.gov/articles/PMC3174607/}{PubMed Central} Choanoflagellate genomes (King et al.~2008, \emph{Nature}; Fairclough et al.~2013) encode neurosecretory SNAREs, complexin, synaptotagmin, voltage-gated Ca²⁺ channels, \href{https://academic.oup.com/mbe/article/31/9/2342/2925718}{Oxford Academic} and a Nav-like channel \href{https://www.pnas.org/doi/10.1073/pnas.1106363108}{PNAS} with a DEEA selectivity filter intermediate between Cav and Nav. \href{https://www.ncbi.nlm.nih.gov/books/NBK207163/}{NCBI} \href{https://www.researchgate.net/figure/Maximum-likelihood-phylogeny-of-the-voltage-gated-sodium-channel-family-The-common_fig5_227857817}{ResearchGate} They encode \textbf{Homer, Shank, and PSD-95/DLG scaffolds}. \href{https://journals.biologists.com/jeb/article/218/4/506/14126/The-origin-and-evolution-of-synaptic-proteins}{The Company of Biologists} \href{https://academic.oup.com/mbe/article/31/9/2342/2925718}{Oxford Academic} Capsasporan filastereans extend the inventory further with Brachyury, Hippo, NF-κB, p53, \href{https://elifesciences.org/articles/77598}{eLife} and a wide GPCR/tyrosine-kinase repertoire (Suga et al.~2013; Sebé-Pedrós et al.~2016, \emph{Cell}). Yañez-Guerra, Thiel and Jékely (2022) found phoenixin and nesfatin-like neuropeptide precursors in choanoflagellates with conserved cleavage sites \href{https://link.springer.com/article/10.1007/s10071-023-01776-z}{Springer} \href{https://www.researchgate.net/publication/378432679_Evolutionary_origin_of_the_nervous_system_from_Ctenophora_prospective}{ResearchGate} --- \textbf{secreted peptide signalling preceded neurons}. \href{https://www.annualreviews.org/content/journals/10.1146/annurev-cellbio-111822-124041?crawler=true}{Annual Reviews}

\textbf{Sponges} lack neurons but possess most of the postsynaptic toolkit: Sakarya et al.~(2007, \emph{PLoS ONE}) showed \emph{Amphimedon queenslandica} encodes \href{https://doaj.org/article/93cf76d6361d47a4b4bde8d2ea666f82}{DOAJ} DLG/PSD-95, Shank, Homer, GKAP, CaMKII and Lin7 \href{https://journals.plos.org/plosone/article?id=10.1371/journal.pone.0000506}{PLOS} with PDZ-binding motifs essentially identical to humans. Conaco et al.~(2019) found that \textasciitilde125 synaptic-gene homologs are co-expressed during sponge metamorphosis \href{https://www.nature.com/articles/s41598-019-51282-x}{Nature} but \textbf{not} as a coordinated synaptic module --- they sit in their original intracellular contexts. Notably absent from sponges: ionotropic glutamate receptors, neurexin/neuroligin, synapsin, RIM, CASK. \href{https://www.ncbi.nlm.nih.gov/pmc/articles/PMC6823388/}{nih}

\textbf{Placozoans} (\emph{Trichoplax adhaerens}) have no neurons, no synapses, and no muscles, but elaborate peptidergic signalling. Varoqueaux et al.~(2018, \emph{Curr Biol}) identified multiple proneuropeptide genes expressed in non-overlapping cell populations across all three layers; eleven peptides applied exogenously elicit reproducible behaviours. \href{https://pubmed.ncbi.nlm.nih.gov/30344118/}{PubMed} Senatore, Reese and Smith (2017) showed peptidergic integration of behaviour despite the absence of synapses. \href{https://elifesciences.org/articles/26349}{eLife} \textbf{Paracrine peptide--GPCR signalling antedates synaptic transmission}, \href{https://www.biorxiv.org/content/10.1101/360925v1}{bioRxiv} and this is the basis of Jékely's (2021, \emph{Phil Trans R Soc B}) \textbf{``chemical brain hypothesis''}: \href{https://www.sciencedirect.com/science/article/abs/pii/S095943882100129X}{ScienceDirect} nervous systems evolved by adding speed and spatial precision to a pre-existing chemical communication network, not by inventing chemical communication.

\textbf{Cnidarians} have unambiguous chemical synapses, often bidirectional with vesicles on both sides of the cleft \href{https://pmc.ncbi.nlm.nih.gov/articles/PMC5285349/}{PubMed Central} (Anderson 1985; Anderson \& Spencer 1989), heavily peptidergic with peptide-gated DEG/ENaC channels \href{https://pmc.ncbi.nlm.nih.gov/articles/PMC11087051/}{PubMed Central} (Assmann et al.~2014). Sebé-Pedrós et al.~(2018, \emph{Cell}) used scRNA-seq in \emph{Nematostella} to show neurons, cnidocytes and gland cells share a common multipotent secretory progenitor --- supporting a \textbf{secretory-cell origin of the neuron}. Thiel et al.~(2024, \emph{eLife}) deorphanised the \emph{Nematostella} peptide GPCR repertoire and showed that cnidarian and bilaterian peptide-receptor pairings have \textbf{no direct orthology} --- peptide signalling expanded independently in the two lineages from a small ancestral set. \href{https://pmc.ncbi.nlm.nih.gov/articles/PMC11087051/}{PubMed Central} \href{https://elifesciences.org/articles/90674}{eLife}

\textbf{Ctenophores} (comb jellies) precipitated the most contested debate in nervous-system evolution. Ryan et al.~(2013, \emph{Science}; \emph{Mnemiopsis} genome) and Moroz et al.~(2014, \emph{Nature}; \emph{Pleurobrachia} genome) reported phylogenomic placement as sister to all other animals, \href{https://www.nature.com/articles/nature13400}{Nature +2} supported now by chromosome-scale synteny evidence (Schultz et al.~2023, \emph{Nature}); \href{https://www.ncbi.nlm.nih.gov/pmc/articles/PMC10232365/}{nih} \href{https://pubmed.ncbi.nlm.nih.gov/37198475/}{PubMed} the absence or non-neuronal expression of canonical neurogenic transcription factors \href{https://pmc.ncbi.nlm.nih.gov/articles/PMC4337882/}{PubMed Central} and many classical neurotransmitter pathway genes; \href{https://www.nature.com/articles/nature13400}{Nature} \href{https://www.researchgate.net/publication/370154444_Syncytial_nerve_net_in_a_ctenophore_adds_insights_on_the_evolution_of_nervous_systems}{ResearchGate} ctenophore iGluRs that bind glycine rather than glutamate; and a partly \textbf{syncytial} subepithelial nerve net documented by volume EM (Burkhardt et al.~2023, \emph{Science}). \href{https://www.biorxiv.org/content/10.1101/2025.06.26.661250v2.full}{bioRxiv +2} Moroz argues this represents independent, parallel origin of neurons. \href{https://pubmed.ncbi.nlm.nih.gov/25696823/}{PubMed} \href{https://www.semanticscholar.org/paper/Convergent-evolution-of-neural-systems-in-Moroz/7ede2e7fe68e14c1b1578a267cc649ce600a9397}{Semantic Scholar} The counter-position, defended by Burkhardt and Jékely (2021) in their ``divide-and-conquer'' model, identifies a Unc13-RIM spiking sensory-peptidergic cell at the \emph{Mnemiopsis} mouth that may be homologous with bilaterian neurons. \href{https://www.sciencedirect.com/science/article/abs/pii/S095943882100129X}{ScienceDirect} The honest reading is that ctenophores possess the pan-eukaryotic substrate for chemical synapses but their neuron complement, neurotransmitter chemistry and nerve-net architecture diverge sufficiently that whether their ``neurons'' are homologous to bilaterian neurons depends on what one defines as the neuron's identity. Recent work (Copley 2025, \emph{MBE}) \href{https://www.ncbi.nlm.nih.gov/pmc/articles/PMC12728501/}{nih} has challenged the synteny argument, and the question is \textbf{not closed}.

The \textbf{modern bilaterian synapse} therefore assembled in stages: pre-metazoan SNARE/SM secretion, Ca²⁺ exocytosis, voltage-gated channels, GPCRs, peptide precursors and PDZ/MAGUK scaffolds in choanoflagellates; full PSD scaffolds and peptidergic networks at the metazoan stem; bona fide chemical synapses with canonical pre/post architecture in the eumetazoan ancestor; the modern presynaptic active zone (RIM, Munc13, ELKS, Bassoon, Piccolo, liprin-α tightly co-organised) and polarised unidirectional synapses at the bilaterian stem; and dramatic vertebrate-specific paralog expansion driven by two whole-genome duplications \href{https://www.nature.com/articles/ncomms14613}{Nature} (Bayés et al.~2011, 2017). \textbf{Electrical synapses} complicate the picture: innexins/pannexins are the ancestral metazoan gap-junction proteins, \href{https://link.springer.com/article/10.1186/s12862-019-1369-4}{Springer} and connexins are a chordate-specific family \href{https://link.springer.com/article/10.1186/s12862-019-1369-4}{Springer} with no sequence homology to innexins. Welzel and Schuster (2022, \emph{eLife}) argue that early chordates lost innexin diversity and connexins replaced them --- a textbook case of one protein family substituting for another in conserved cellular function.

\subsection{What the urbilaterian probably had}\label{04_urbilateria-what-the-urbilaterian-probably-had}

Reconstructing the last common bilaterian ancestor --- the urbilaterian, dating to roughly 600--650 Ma based on relaxed molecular clocks (Peterson et al.~2008; Erwin et al.~2011; dos Reis et al.~2015, \emph{Curr Biol}) --- requires triangulation among living phyla and is genuinely contested. Two camps persist (Northcutt 2012, \emph{PNAS}). \href{https://karger.com/bbe/article/82/4/211/326147/How-Basal-Are-the-Basal-Ganglia-Commentary-on}{Karger Publishers}

The \textbf{complex CNS camp} --- Arendt, De Robertis, Reichert, Hirth, Strausfeld --- points to conserved anteroposterior (Six3 → Otx → Gbx → Pax2/5/8 → Hox) and mediolateral (Nk2.2 → Nk6 → Pax6 → Msx) patterning across vertebrates, amphioxus, ascidians, \emph{Drosophila}, \emph{Platynereis} and even hemichordate ectoderm. Denes et al.~(2007, \emph{Cell}) provided the strongest single-paper case: the \emph{Platynereis dumerilii} nerve cord shows mediolateral neurogenic columns matching the vertebrate neural tube gene-for-gene. Hirth et al.~(2003) and Lichtneckert and Reichert (2005, \emph{Heredity}) argued that Otx, Pax2/5/8 and Hox demarcate three brain regions in both \emph{Drosophila} and vertebrates, suggesting a tripartite urbilaterian brain. Tomer, Denes, Tessmar-Raible and Arendt (2010, \emph{Cell}) used image registration to argue that annelid mushroom bodies and the vertebrate pallium develop from corresponding molecular coordinates --- a radical claim of deep homology of higher integrative centres. Tessmar-Raible et al.~(2007) found conserved sensory-neurosecretory cell types in annelid and fish forebrain, supporting hypothalamic homology. Arendt, Tosches and Marlow (2016) propose the bilaterian brain emerged from fusion of an ancestral apical (Six3+, neurosecretory, derived from the cnidarian aboral pole organ) and blastoporal (Hox+, motor, around the mouth) nervous system.

The \textbf{simpler urbilaterian camp} --- N. D. Holland, Lowe, Gerhart, Kirschner, Hejnol, Martindale, Martín-Durán --- emphasises that hemichordate \emph{Saccoglossus} has a diffuse epidermal nervous system despite chordate-like AP ectodermal patterning (Lowe et al.~2003, \emph{Cell}; 2006, \emph{PLoS Biol}); that hemichordate BMP signalling does not segregate neural from epidermal ectoderm; that xenacoelomorphs ancestrally have a basiepidermal nerve net without nephrozoan-style CNS architecture (Cannon et al.~2016, \emph{Nature}); and that \textbf{dorsoventral nerve-cord patterning is not conserved across Bilateria}. Martín-Durán et al.~(2018, \emph{Nature}) examined Xenacoelomorpha, Rotifera, Nemertea, Brachiopoda and Annelida and found no conserved D-V molecular regionalisation --- even \emph{Owenia fusiformis}, a basal annelid with vertebrate-like trunk neuroanatomy, lacks the staggered Nk-Pax pattern. They argue convergent evolution of medially condensed nerve cords is more parsimonious than a single ancestral CNS. On this view, the urbilaterian had the \textbf{genetic toolkit} for a complex CNS but expressed it as a regionalised, basiepidermal nerve plexus; centralised cords evolved convergently in chordate, arthropod-like and annelid lineages.

The dorsoventral inversion hypothesis --- Geoffroy Saint-Hilaire's 1822 observation that arthropod and vertebrate body plans look like inverses, revived by Arendt and Nübler-Jung (1994) and given molecular force by De Robertis and Sasai (1996, \emph{Nature}) --- sits at the heart of this debate. The BMP-Chordin and Dpp-Sog antagonist pairs are homologous but mirror-inverted in expression; the CNS forms on the low-BMP side in both phyla. Holley et al.~(1995) showed Sog rescues Chordin function and vice versa. But Lowe et al.~(2006) showed hemichordates retain the BMP-Chordin axis without using it to segregate neural from epidermal ectoderm. Whether literal 180° inversion occurred or whether dorsalised and ventralised CNS condensations evolved convergently from a diffuse ancestor remains genuinely open.

What is reasonably established: the BMP-Chordin axis is bilaterian-ancestral; Six3, Otx, Gbx and Hox demarcate broad axial territories shared between cnidarians (oral-aboral) and bilaterians (anteroposterior), so the \textbf{axial regulatory program is pre-bilaterian}; Thiel et al.~(2018) showed that more than 20 neuropeptide GPCR families already existed at the bilaterian stem, \textbf{predating} the formation of complex CNS. The urbilaterian had at least the genetic toolkit for nervous system regionalisation, peptidergic and small-molecule chemical signalling, and probably some integration centres --- apical neurosecretory and blastoporal/feeding-locomotor. Whether it had a morphologically condensed CNS, whether D-V inversion occurred, and whether tripartite brain organisation is homologous all remain contested in 2025--2026.

\subsection{Conservation and convergence: lessons from invertebrate models}\label{04_urbilateria-conservation-and-convergence-lessons-from-invertebrate-models}

A handful of invertebrate models illustrate how much molecular substrate is conserved across bilaterians and how often circuit architecture converges on similar solutions. They are not the trunk of the human lineage, but they discipline inference about the urbilaterian condition and illustrate the through-line of the chapter.

\textbf{C. elegans} has exactly 302 neurons in the hermaphrodite, 385 in the male, with complete EM connectomes for both sexes (White et al.~1986; Cook et al.~2019, \emph{Nature}). The worm uses essentially the full bilaterian neurotransmitter palette, with dopamine produced in just four pairs of neurons. The central conceptual lesson --- articulated forcefully by Bargmann (2012, \emph{BioEssays}), Bargmann and Marder (2013, \emph{Nat Methods}), and Marder (2012, \emph{Neuron}) --- is that \textbf{the wiring diagram is necessary but not sufficient}: the same anatomy generates many functional circuits depending on neuromodulatory state. Flavell et al.~(2013, \emph{Cell}) showed that mutually inhibitory serotonin and PDF neuropeptide systems toggle persistent roaming versus dwelling foraging states using identical underlying circuitry. The implication for any account of agency is that ``the connectome'' cannot be the whole story; molecules in transit are part of the substrate. Crucially, \emph{C. elegans} is a highly derived ecdysozoan, and its tiny neuron count is secondary miniaturisation, not the ancestral state.

\textbf{Drosophila} mushroom bodies foreshadow themes the chapter saves for later vertebrate treatment but that already operate in invertebrate brains. The roughly 130 dopaminergic neurons of the PPL1 and PAM clusters provide compartmentalised modulatory input to Kenyon cells, gating long-term depression at KC→MBON synapses to bias approach versus avoidance (Aso et al.~2014, \emph{eLife}; Cognigni, Felsenberg \& Waddell 2018; Yamagata et al.~2015, \emph{PNAS}). PPL1 dopaminergic neurons signal punishment; PAM neurons convey reward --- a logic strikingly parallel to vertebrate basal-ganglia/dopamine reward-prediction-error systems. Whether mushroom body and pallium are homologous (Tomer et al.~2010) or convergent (Northcutt 2012; Loesel commentaries) is one of the enduring debates in evo-neuro: \textbf{shared transcription-factor codes can reflect regulatory cooption rather than structural homology}, and few-phylum sampling makes parsimony assignments unstable.

\textbf{Planarians} demonstrate a minimal bilaterian centralised brain --- bilobed cephalic ganglia connected to two ventral nerve cords --- with conserved bilaterian neurotransmitter and neuropeptide systems and the capacity for complete brain regeneration in seven days from pluripotent neoblasts (Reddien \& Sánchez Alvarado 2004; Wagner, Wang \& Reddien 2011, \emph{Science}). They illustrate that neural cell-type identity is decoupled from positional history.

\textbf{\emph{Platynereis dumerilii}} has been Arendt's lab's chosen ``slow-evolving'' model for the urbilaterian. Its trochophore larva and adult brain have provided the strongest evidence --- contested though it is --- for deep bilaterian homologies: ciliary photoreceptors with vertebrate-type c-opsin in the larval brain alongside rhabdomeric photoreceptors with r-opsin in the eyespots (Arendt et al.~2004, \emph{Science}), implying both photoreceptor lineages coexisted in the urbilaterian; mediolateral nerve-cord patterning matching the vertebrate neural tube (Denes et al.~2007); and the proposed mushroom-body/pallium homology (Tomer et al.~2010). Verasztó et al.~(2024, \emph{Nat Commun}) provides the larval connectome with 202 neuronal cell types. Whether one accepts the strong homology readings or treats them as evidence of conserved gene-regulatory networks deployed convergently, \emph{Platynereis} shows that the molecular fingerprints of vertebrate neuron classes have deep bilaterian roots.

The cross-cutting theme is \textbf{same molecules, different circuits}: glutamate, GABA, ACh, monoamines and neuropeptides are universal bilaterian currencies, but the wiring topologies that deploy them have diverged enormously. Arendt (2008, \emph{Nat Rev Genet}) introduced the \textbf{molecular fingerprint} concept of cell-type identity, and the cell-type-homology versus organ-homology distinction (Arendt et al.~2016, \emph{Nat Rev Genet}) has become essential vocabulary for thinking about what is and is not preserved across 600 million years.

\subsection{The deuterostome turn}\label{04_urbilateria-the-deuterostome-turn}

Modern phylogenomics consolidates Deuterostomia as \textbf{Ambulacraria (echinoderms + hemichordates) + Chordata} (Bourlat et al.~2006, \emph{Nature}; Telford et al.~2015; Cannon et al.~2016, \emph{Nature}). Crown deuterostomes date to roughly 573--557 Ma, in the late Ediacaran, coincident with the Shuram excursion and after the Gaskiers glaciation. The phylogenetic placement of Xenacoelomorpha --- sister to Nephrozoa (Cannon et al.~2016) versus sister to Ambulacraria within ``Xenambulacraria'' (Philippe et al.~2019, \emph{Curr Biol}) --- remains unresolved and matters for inference about whether the deuterostome ancestor was simple or complex.

\textbf{Hemichordates} are critical because they break the assumption that AP regulatory conservation entails CNS centralisation. Adult enteropneusts have a basiepithelial nerve plexus concentrated in the proboscis plus paired dorsal and ventral nerve cords; the dorsal collar cord is internalised by a neurulation-like process (Kaul-Strehlow \& Stach 2010). Lowe et al.~(2003, \emph{Cell}) showed that 22 orthologs of chordate AP-patterning genes --- \emph{Otx}, \emph{En}, \emph{Gbx}, \emph{Hox1--13}, \emph{Nkx2.1}, \emph{Pax6}, \emph{Bf-1}, \emph{Dlx}, \emph{Emx} --- are expressed in nearly identical AP order in \emph{Saccoglossus kowalevskii} ectoderm, but as \textbf{circumferential rings}, not in a condensed CNS. Pani et al.~(2012, \emph{Nature}) found anterior signalling centres (Six3, Fox, Rx) homologous to vertebrate forebrain organisers. Andrade López et al.~(2023, \emph{PLoS Biol}) provided a comprehensive molecular CNS map of \emph{S. kowalevskii}, finding that the proboscis-base plexus is the major integrative centre, the dorsal and ventral cords are largely through-conduction tracts, and \textbf{no clean CNS/PNS division} exists. The hemichordate evidence therefore argues that the AP gene-regulatory network is dissociable from morphological centralisation: the deuterostome ancestor may have had a regionalised diffuse nervous system, with chordate-style centralisation as a derived event.

\textbf{Echinoderms} are derived bilaterians whose dipleurula-type larvae are bilaterally symmetric with apical organs; pentaradiality emerges at metamorphosis from the left coelom. The adult CNS is a circumoral nerve ring plus five radial cords, with no centralised brain. The fossil record (Zamora et al.~2012; Zamora \& Rahman 2014) documents stepwise acquisition through bilateral \emph{Ctenoimbricata} and ctenocystoids, asymmetric cinctans, triradiate helicoplacoids, and pentaradial helicocystoids. Formery et al.~(2023, \emph{Nature}) showed that the bilaterian AP axis maps onto the medio-lateral axis of each ray in asteroid juveniles. Pentaradiality is a \textbf{derived loss of cephalisation} layered on a preserved bilaterian patterning toolkit --- a useful counterexample to any narrative of inevitable cephalic progress.

The \textbf{chordate body plan} assembled five diagnostic features: the notochord (a dorsal mesodermal rod induced by Brachyury and secreting Hh signals to specify floor plate); the dorsal hollow nerve cord, formed by neurulation under a BMP-Chordin dorsoventral gradient and SHH from notochord and floor plate; pharyngeal slits or gill bars (already present in hemichordates); the endostyle, a ventral pharyngeal groove and thyroid homolog expressing Nkx2.1 and TPO; and the post-anal tail with segmented somites in vertebrates and amphioxus. Chordates date to roughly 550 Ma by molecular clock, with first unambiguous fossils --- \emph{Cathaymyrus} (Shu et al.~1996), \emph{Pikaia} (Conway Morris \& Caron 2012, \emph{Biol Rev}), and the early stem vertebrates \emph{Myllokunmingia} and \emph{Haikouichthys} (Shu et al.~1999, 2003, \emph{Nature}) --- appearing in the Chengjiang biota at \textasciitilde518 Ma.

Why a dorsal nerve cord? The Geoffroy/De Robertis dorsoventral inversion remains the standard answer, with the qualifier that hemichordate evidence is mixed: the hemichordate dorsal collar cord is on the same side as the chordate neural tube relative to the mouth, arguing against simple inversion or supporting a pre-chordate inversion (Nübler-Jung \& Arendt 1999). Yu et al.~(2007, \emph{Nature}) showed that amphioxus BMP/Chordin behaves like vertebrates. Luo et al.~(2019, \emph{PNAS}) showed that BMP manipulation in indirect-developing hemichordates produces phenotypes resembling a proposed chordate intermediate. The mechanistic story is therefore that \textbf{the BMP-Chordin cassette was tinkered with} to produce the dorsal hollow tube, perhaps from a pelagic indirect-developing chordate ancestor.

\subsection{Amphioxus, tunicates and the chordate brain}\label{04_urbilateria-amphioxus-tunicates-and-the-chordate-brain}

Phylogenomics produced one of the great recent reorganisations of the chordate tree. Delsuc, Brinkmann, Chourrout and Philippe (2006, \emph{Nature}), using 146 nuclear genes, overturned the classical Euchordata hypothesis: \textbf{tunicates are sister to vertebrates} (the clade Olfactores), not amphioxus. Cephalochordates --- amphioxus --- are the most basally branching extant chordate. This has been confirmed repeatedly (Bourlat et al.~2006; Putnam et al.~2008; Delsuc et al.~2018, \emph{BMC Biol}).

But the chordate tree carries a ``tunicate problem.'' Tunicates evolve roughly twice as fast as other chordates (Tsagkogeorga et al.~2010), with massive genome rearrangement and gene loss; \emph{Oikopleura dioica} is hyperderived (Denoeud et al.~2010, \emph{Science}). \textbf{Tunicate ``simplicity'' is largely secondary}, not ancestral retention. Despite branching earlier, \textbf{amphioxus is the better proxy for the ancestral chordate brain plan} --- slower-evolving, with a single Hox cluster (versus the four vertebrate clusters from two rounds of whole-genome duplication), a retained chordate body plan, and roughly 20,000 neurons in its CNS (Candiani et al.~2012; Putnam et al.~2008, \emph{Nature}).

The amphioxus brain has no morphologically distinct fore-, mid-, and hindbrain, only a \textbf{cerebral vesicle} at the anterior tip of the dorsal nerve cord. But molecular regionalisation is already present (Albuixech-Crespo et al.~2017, \emph{PLoS Biol}; Castro et al.~2006; Holland 2015): \emph{Otx} anterior, \emph{Gbx} abutting \emph{Otx} posteriorly to form a molecular MHB-equivalent boundary (lacking, however, the \emph{En} + \emph{Pax2/5/8} + \emph{Wnt1} co-expression that defines a true vertebrate MHB organiser); nested \emph{Hox} domains in the hindbrain-equivalent suggesting rhombomere-like molecular subdivisions without overt segmentation; \emph{FoxG1}, \emph{Emx}, \emph{Lhx2/9} marking a putative telencephalic-like region; \emph{Nkx2.1} in ventral diencephalic region and endostyle; \emph{Pax6} in frontal eye Row 1 cells. The amphioxus photoreceptive organs --- frontal eye, lamellar body (a pineal homolog candidate), Joseph cells, organs of Hesse, and the infundibular organ secreting Reissner's fibre --- are recognisable in vertebrate terms (Lacalli 1996, 2018; Vopalensky et al.~2012, \emph{PNAS}). The ancestral chordate brain, on this evidence, was already regionalised with molecular homologs of vertebrate diencephalon and posterior brain, but lacked a true MHB organiser and rhombomeric segmentation.

\textbf{Ciona intestinalis} larvae provide the smallest fully reconstructed chordate connectome: 177 neurons in the CNS, 6,618 chemical synapses, 1,772 neuromuscular junctions, and 1,206 gap junctions (Ryan, Lu \& Meinertzhagen 2016, \emph{eLife}). The brain vesicle, sensory ocellus and otolith, neck, motor ganglion, and caudal nerve cord express conserved chordate markers --- \emph{Otx} anterior, \emph{Pax2/5/8} with \emph{En} at MHB-like position, \emph{Hox} posterior --- with a distinctive marked left-right asymmetry. The ascidian adenohypophysis homolog forms from an anterior placode (Boorman \& Shimeld 2002). At metamorphosis the sessile adult ascidian regresses its nervous system to a cerebral ganglion. Whether the ancestral chordate was pelagic (Garstang's auricularia/paedomorphosis hypothesis) or sessile (Berrill) was debated for a century; the modern view, supported by Luo et al.~(2019), is that the chordate ancestor was a free-swimming animal, and ascidian sessility is secondary.

The \textbf{chordate-specific innovations} beyond the ancestral bilaterian condition were therefore: the dorsal hollow nerve tube formed by neurulation under BMP-Chordin and Brachyury/Hh control; the notochord as inducer; the endostyle; segmented somites; and the further elaboration in Olfactores of an MHB organiser. Vertebrates added on top of this neural crest, ectodermal placodes producing paired sensory organs, and two rounds of whole-genome duplication (Dehal \& Boore 2005) that produced the paralog-rich vertebrate brain proteome.

\subsection{The lamprey window}\label{04_urbilateria-the-lamprey-window}

Cyclostomes --- lampreys and hagfishes --- are the only extant jawless vertebrates, and Heimberg et al.~(2010, \emph{PNAS}) established their monophyly using shared microRNA families, paralogues, and substitutions, with reanalysis of morphology revealing that prior support for paraphyly rested on character-coding errors. The agnathan-gnathostome split is placed at roughly 500 Ma, in the Cambrian-Ordovician (Kuraku 2008; Donoghue \& Smith 2004). Lampreys cannot be treated as frozen ancestors --- they have had half a billion years of independent evolution including their own genome duplications, parasitic adult feeding, programmed somatic DNA elimination during early development (Smith et al.~2018, \emph{Nat Genet}), and the dramatic ammocoete metamorphosis. But triangulating among lamprey, hagfish, and stem-group fossils makes them an unrivalled window onto the ancestral vertebrate condition.

The \textbf{earliest unambiguous vertebrates} appear in the Lower Cambrian Chengjiang biota at \textasciitilde518--520 Ma. \emph{Haikouichthys} and \emph{Myllokunmingia} (Shu et al.~1999, 2003, \emph{Nature}) preserve lobate heads with paired eyes, possible nasal sacs and otic capsules, branchial arches, vertebral elements, and zigzag myomeres --- strongly suggesting that the \textbf{tripartite brain plan, neural-crest-derived sensory placodes, and paired camera eyes were present in stem vertebrates} (Northcutt 2002). Tian et al.~(2022, \emph{Science}) reinterpreted yunnanozoans as stem-vertebrates with cellular cartilage. Conodonts and ostracoderms across the Ordovician-Devonian are now read as \textbf{stem gnathostomes}, not stem cyclostomes (Janvier 2008, 2015), and Gai et al.~(2011, \emph{Nature}) showed that the separation of paired olfactory organs from the hypophyseal duct occurred on the gnathostome stem after cyclostomes diverged. \textbf{The basic vertebrate brain plan was therefore stabilised by \textasciitilde520 Ma, before any major living lineage diverged.}

The adult lamprey brain shows the vertebrate-typical tripartite plan (Pombal \& Megías 2017; Suryanarayana et al.~2022, \emph{Brain Behav Evol}). The \textbf{telencephalon} has paired olfactory bulbs and hemispheres; the dorsal pallium, contrary to older views that it was largely olfactory, contains discrete motor, somatosensory and visual areas with retinotopic mapping (Suryanarayana et al.~2017, 2020). The subpallium contains striatum, pallidum (with GPi/GPe homologs), septum, and a habenula-projecting pallidum. Both LGE-derived and MGE-derived inhibitory progenitor zones are present (Lamanna et al.~2023, \emph{Nat Ecol Evol}). The \textbf{diencephalon} has a dorsal thalamus with relay and higher-order nuclei projecting to pallium and striatum, prethalamus, hypothalamus, posterior tuberculum (the lamprey homolog of the gnathostome SNc/A11 dopaminergic centre), and an epithalamus including pineal and parapineal organs and an asymmetric habenula (Stephenson-Jones et al.~2012, \emph{PNAS}). The \textbf{mesencephalon} contains a well-developed optic tectum (homolog of the superior colliculus) with retinotopic visual map. The \textbf{rhombencephalon} contains a large reticular formation with the giant Müller cells and the Mauthner cell driving reticulospinal commands; an octavolateral area serving electroreception, mechanoreception and vestibular input; and branchiomotor and somatomotor nuclei in rhombomeric register (Murakami et al.~2004).

What is \textbf{already vertebrate-like} in lamprey and inferred ancestral: neural crest and ectodermal placodes, paired camera eyes with lens, pineal/parapineal photoreceptors, olfactory placode and bulb, lateral-line and electroreceptive systems, octavolateral nuclei, rhombomeric Hox patterning, basal-ganglia direct/indirect pathway architecture, habenula bipartite organisation with asymmetry, hypothalamic-pituitary axis with neuropeptide control, and Schwann-cell and astrocyte orthologs.

What is \textbf{primitive or absent}: jaws, paired appendages with endoskeletal girdles, true compact myelin, oligodendrocytes, cellular bone, Ig/TCR/RAG-based adaptive immunity (lampreys use somatically diversified Variable Lymphocyte Receptors built from leucine-rich repeats --- Pancer et al.~2004, \emph{Nature}; Boehm et al.~2012), sympathetic ganglia, and a true cerebellum. The myelin question deserves care: Bullock et al.~(1984) showed lampreys lack myelin; Smith et al.~(2013) reported a fragment annotated as MBP, but Werner (2013) showed it is actually GOLLI; Weil et al.~(2018, \emph{J Neurosci}) found that lamprey peripheral axons are ensheathed by Schwann-like cells expressing SoxE2/E3 orthologs and CNS axons by astrocyte-like cells, indicating axon ensheathment and glial metabolic support predate myelin but \textbf{compact myelin itself is a gnathostome innovation}. The cerebellum case is similar: Lamanna et al.~(2023) found no granule or Purkinje cell signatures in adult lamprey single-cell data, although Sugahara et al.~(2016, \emph{Nature}) reported transient Purkinje- and granule-cell transcription factors in dorsal rhombomere 1 of embryos --- a true cerebellum with rhombic-lip-derived granule cells and inferior olive evolved on the gnathostome stem; cerebellum-like circuits in octavolateralis nuclei are older.

The \textbf{neurotransmitter inventory} of the lamprey is the central evidence that the chemical architecture of vertebrate behaviour was already in place before the gnathostome divergence. \textbf{Dopamine} is most thoroughly characterised. Dopaminergic neurons map closely to the gnathostome A8--A16 scheme (Pombal et al.~1997; Barreiro-Iglesias et al.~2008; Pérez-Fernández et al.~2014, \emph{J Neurosci}): olfactory bulb (A16-equivalent), hypothalamus and preoptic area (A11--A15), posterior tuberculum (the lamprey SNc/VTA homolog projecting to striatum, tectum and habenula), plus retinal, spinal and CSF-contacting populations. D1 and D2 receptors are present and pharmacologically conventional; Robertson et al.~(2012) cloned the D2, and Pérez-Fernández, Megías and Pombal (2016, \emph{Front Neuroanat}) the D4. \textbf{Dopaminergic input to striatum modulates direct and indirect pathways exactly as in mammals} (Ericsson et al.~2013; Pérez-Fernández et al.~2014). \textbf{Serotonin} populates raphe-equivalent isthmic and rhombencephalic nuclei plus pineal, hypothalamus and spinal CSF-contacting cells (Barreiro-Iglesias et al.~2008, \emph{J Chem Neuroanat}); spinal serotonin modulates swimming frequency and slow afterhyperpolarisation in motor-network neurons (Wallén et al.~1989). \textbf{GABA and glutamate} are ubiquitous; lamprey striatal projection neurons are GABAergic with substance-P (D1-direct) and enkephalin (D2-indirect) subtypes, exactly as in mammals (Stephenson-Jones et al.~2011, \emph{Curr Biol}). Co-localisation of dopamine with GABA and serotonin with GABA in subsets of lamprey neurons (Barreiro-Iglesias et al.~2009) suggests neurotransmitter co-release predates jawed vertebrates. \textbf{Acetylcholine} marks cranial and spinal motor neurons and mesopontine nuclei homologous to the pedunculopontine and laterodorsal tegmental, which connect to all basal-ganglia nuclei (Stephenson-Jones et al.~2012, \emph{J Comp Neurol}). \textbf{Noradrenaline} is genuinely contested: Barreiro-Iglesias et al.~(2010) found no clear locus coeruleus homolog by tyrosine hydroxylase ISH and immunohistochemistry, and sympathetic ganglia are absent. Current consensus places a definitive locus coeruleus and the noradrenergic projection system on the \textbf{gnathostome stem}. \textbf{Histamine, GnRH (multiple forms), vasotocin/oxytocin-like peptides, NPY, CCK, substance P, enkephalin, galanin, somatostatin and CRH} have all been described with conserved distributions; the hypothalamic-pituitary axis with adenohypophyseal releasing factors is established in agnathans (Sower 2018).

The single most consequential finding is the Karolinska group's demonstration that \textbf{all canonical basal-ganglia components} are present in lamprey (Stephenson-Jones et al.~2011, 2012; Robertson et al.~2014; Grillner \& Robertson 2016, \emph{Curr Biol}; Suryanarayana et al.~2022): dorsal striatum with D1-substance P and D2-enkephalin spiny projection neurons, GPi/GPe equivalents, glutamatergic STN, dopaminergic SNc-equivalent in the posterior tuberculum, and a habenula-projecting pallidum. Direct and indirect pathways, transmitter phenotypes, ion-channel complement (Kir), and tonic GABAergic output to brainstem motor centres are all conserved. \textbf{The action-selection architecture predates gnathostomes by at least 500 million years.} Detailed treatment is reserved for later chapters; the load-bearing point here is that the basal-ganglia solution to action selection --- a vertebrate elaboration on top of inherited bilaterian peptidergic and small-molecule signalling --- was already in place when our lineage diverged from cyclostomes.

Lamprey swimming has been the gold-standard vertebrate central pattern generator since the 1980s (Grillner et al.~1991, \emph{Annu Rev Neurosci}; Grillner 2003, 2021). Spinal segmental networks of glutamatergic excitatory interneurons, glycinergic commissural inhibitory interneurons, motor neurons and sensory edge cells generate alternating left-right bursts; intersegmental coupling yields a travelling wave; NMDA-receptor-driven plateau properties shape low-frequency bursts; reticulospinal command neurons in mesencephalic and diencephalic locomotor regions initiate swimming; descending dopamine, serotonin and GABA modulate rhythm and intensity. \textbf{The same circuit logic operates in mouse, zebrafish and \emph{Xenopus} tadpole.}

Hagfish --- even more reduced eyes (lensless), no electroreception, no dorsal octavolateral nucleus, simpler telencephalon --- retain the basic vertebrate tripartite brain plan and many neural-crest derivatives. Some ``primitive'' features of cyclostomes are secondary losses; the ancestral-vertebrate inference requires triangulation among lamprey, hagfish and stem-group fossils.

\subsection{What can break, and what was added later}\label{04_urbilateria-what-can-break-and-what-was-added-later}

The chapter's through-line is that agency has a material substrate that can fail at every level. The same molecular fragilities that produce human disease are recognisable across 600 million years of bilaterian history. \textbf{Tyrosine hydroxylase loss} produces dopamine deficits and impaired learning in flies and movement in worms, and is the proximate cause of Parkinsonism in humans. \textbf{GAD loss} produces hyperexcitability across nematodes, flies and mammals. \textbf{Vesicular transporter (VGAT, VAChT, VMAT) mutations} produce transmitter release failure. \textbf{Channelopathies} in Na⁺, K⁺ or Ca²⁺ channels --- \emph{unc-2} and \emph{egl-19} in \emph{C. elegans}, \emph{Shaker} in \emph{Drosophila}, SCN1A/GABRG2 in human epilepsies --- produce convergent epileptiform or paralytic phenotypes. The substrate is the same, the failures rhyme.

What was added between the urbilaterian and the lamprey was, in concrete terms, surprisingly little chemistry and a great deal of architecture. The neurotransmitter palette was already present at the bilaterian stem; the synaptic protein toolkit was already present at the metazoan stem; the AP and DV regulatory programs were already present in the urbilaterian. What chordates added was the dorsal hollow nerve tube and notochord-based induction; what Olfactores added was the MHB organiser and elaborated rhombomeres; what vertebrates added was neural crest, ectodermal placodes producing paired sensory organs, two rounds of whole-genome duplication producing paralog-rich postsynaptic densities, and --- already in lamprey --- the basal-ganglia direct/indirect pathway architecture for action selection. \textbf{The clearest gnathostome innovations beyond cyclostomes are compact myelin and oligodendrocytes, the true rhombic-lip-derived cerebellum, the locus coeruleus and noradrenergic projection system, the sympathetic chain, and the elaboration of cortical and pallial architecture.}

\subsection{Conclusion}\label{04_urbilateria-conclusion}

Three commitments shape the chapter that follows. First, \textbf{neurotransmitters are not the magic bullets of folk neuroscience}: dopamine is not ``the reward molecule,'' serotonin is not ``the happiness molecule,'' and the temptation to read these molecules as purpose-built signals misrepresents their evolutionary origin. They are tyrosine and tryptophan hydroxylation products co-opted from antioxidant pathways, glutamate decarboxylation products from acid-stress shunts, choline acetylation products from universal metabolism. They were available, and so they were used. Second, \textbf{the synapse is a polarised re-organisation of pre-existing parts}: the SNARE/SM secretion module, the Ca²⁺-triggered exocytosis machinery, the PDZ/MAGUK scaffolds, the GPCRs, the Cys-loop and ionotropic glutamate channels --- every structural component existed in unicellular eukaryotes or early metazoans before the bilaterian synapse assembled. Third, \textbf{agency, in any sense that survives evolutionary scrutiny, is the product of conserved molecular substrate deployed in evolved circuit architecture}, and that architecture solves a specific problem --- action selection in a directionally moving body --- that bilateral symmetry and through-gut feeding posed roughly 600 million years ago.

The lamprey is the chapter's hinge. Its brain already contains the small-molecule and peptidergic signalling repertoires, the basal-ganglia direct and indirect pathways, the habenular asymmetry, the hypothalamic-pituitary axis, the rhombomeric hindbrain organisation, and the action-selection logic that later vertebrate evolution would elaborate but not invent. Where lamprey, hagfish and basal gnathostomes agree, ancestral vertebrate character status is robust. Where they disagree --- cerebellum, electroreception, myelin, locus coeruleus --- parsimony is genuinely contested. The honest reading is that the substrate of agency in our lineage was largely complete by \textasciitilde500 Ma, and the task of subsequent chapters is not to find new chemistry but to trace how circuit architecture elaborated, broke, and was repaired across the half billion years since.

The chapter that follows this one will dwell on the basal ganglia and the dopaminergic mid-brain projections that lamprey already possesses, on the cerebellum that gnathostomes added, on the cortical elaborations of mammals, and on the human-specific tinkering with timing, learning and prediction. None of those later elaborations introduces new chemistry. They reorganise inherited substrate. The substrate is what can break, and tracing that substrate from urbilaterian to early chordate to lamprey is --- materially, biochemically, and historically --- what it means to take the evolutionary view of agency seriously.
